\section{Тест производительности}
Для оценки эффективности разработанного алгоритма поразрядной сортировки было проведено сравнение
со стандартной устойчивой сортировкой std::stable\_sort из стандартной библиотеки C++. В качестве
исходных данных использовался набор из 1 миллиона записей, каждая из которых содержала дату в формате
DD.MM.YYYY и произвольную строку-значение. Тестовые данные были сгенерированы с помощью программы
\texttt{make\_test}, после чего проводился замер времени выполнения сортировки с использованием утилиты
\texttt{time}.

\begin{alltt}
root@66a8fcab42ba:/workspaces/discran-labs/laba-1/build# make
[ 37%] Built target my_sort
[ 75%] Built target std_sort
[100%] Built target make_test
root@66a8fcab42ba:/workspaces/discran-labs/laba-1/build# ./make_test 1000000 > test.txt
root@66a8fcab42ba:/workspaces/discran-labs/laba-1/build# wc -l test.txt
1000000 test.txt
root@66a8fcab42ba:/workspaces/discran-labs/laba-1/build# time ./my_sort < test.txt > my_sort_res.txt

real    0m7.397s
user    0m4.231s
sys     0m1.144s
root@66a8fcab42ba:/workspaces/discran-labs/laba-1/build# time ./std_sort < test.txt > std_sort_res.txt

real    0m8.968s
user    0m6.371s
sys     0m0.486s
root@66a8fcab42ba:/workspaces/discran-labs/laba-1/build# 
\end{alltt}

Как видно из результатов, разработанная реализация поразрядной сортировки показала лучшие результаты:
полное время выполнения составило 7.397 секунды, тогда как \texttt{std::stable\_sort} потребовалось 8.968 секунды,
что примерно на 21\% медленнее.

Преимущество поразрядной сортировки объясняется идеей алгоритма и спецификой входных данных. В отличие от алгоритмов
сравнения, таких как \texttt{std::stable\_sort} (реализована на базе сортировки слиянием), которые имеют теоретическую
нижнюю границу сложности $O(n log n)$, поразрядная сортировка обладает линейной сложностью $O(n * (k + b))$, где $k$ --- количество
разрядов в ключе, $b$ --- размерность системы счисления. В данном случае ключом является дата, преобразованная в строку формата YYYYMMDD,
что дает фиксированное количество разрядов --- 8, а размерность системы счисления равна 10. Таким образом, при достаточно большом
количестве элементов (1 миллион) поразрядная сортировка становится эффективнее, так как количество операций растет линейно, а не логарифмически.

Таким образом, разработанная реализация поразрядной сортировки демонстрирует превосходство над стандартной устойчивой
сортировкой благодаря линейной сложности алгоритма.

\pagebreak

